\documentclass[11pt, a4paper]{article}
% ── Shared preamble for all RTOS lab documents ────────────────────────────────
% Usage: % ── Shared preamble for all RTOS lab documents ────────────────────────────────
% Usage: % ── Shared preamble for all RTOS lab documents ────────────────────────────────
% Usage: \input{../preamble}  (from each labNN/ directory)
% ──────────────────────────────────────────────────────────────────────────────

% ── Packages ──────────────────────────────────────────────────────────────────
\usepackage[a4paper, top=2.5cm, bottom=2.5cm, left=2.5cm, right=2.5cm, headheight=28pt]{geometry}
\usepackage[utf8]{inputenc}
\usepackage[T1]{fontenc}
\usepackage{lmodern}
\usepackage{booktabs}
\usepackage{array}
\usepackage{tabularx}
\usepackage{longtable}
\usepackage{multirow}
\usepackage{colortbl}
\usepackage{xcolor}
\usepackage{titlesec}
\usepackage{enumitem}
\usepackage{hyperref}
\usepackage{fancyhdr}
\usepackage{parskip}
\usepackage{listings}
\usepackage[most]{tcolorbox}
\usepackage{fontawesome5}
\usepackage{microtype}
\usepackage{graphicx}
\usepackage{amsmath}

% ── Colors (KMUTNB brand) ──────────────────────────────────────────────────────
\definecolor{kmutnbBlue}{RGB}{0,56,116}
\definecolor{kmutnbGold}{RGB}{186,146,36}
\definecolor{lightGray}{RGB}{245,245,245}
\definecolor{midGray}{RGB}{200,200,200}
\definecolor{codeGray}{RGB}{248,248,248}
\definecolor{codeFrame}{RGB}{220,220,220}
\definecolor{tipteal}{RGB}{0,105,92}
\definecolor{warnAmber}{RGB}{121,85,0}
\definecolor{checkGreen}{RGB}{27,94,32}

% ── Section formatting ─────────────────────────────────────────────────────────
\titleformat{\section}
  {\large\bfseries\color{kmutnbBlue}}
  {\thesection.}{0.5em}{}
  [\color{kmutnbGold}\titlerule]

\titleformat{\subsection}
  {\normalsize\bfseries\color{kmutnbBlue}}
  {\thesubsection.}{0.5em}{}

\titleformat{\subsubsection}
  {\small\bfseries\color{kmutnbBlue}}
  {\thesubsubsection.}{0.5em}{}

\titlespacing*{\section}{0pt}{1.4ex plus .2ex}{0.6ex plus .1ex}
\titlespacing*{\subsection}{0pt}{1.0ex plus .2ex}{0.4ex plus .1ex}

% ── Listings (C and bash) ──────────────────────────────────────────────────────
\lstdefinestyle{cstyle}{
  language        = C,
  basicstyle      = \ttfamily\small,
  keywordstyle    = \bfseries\color{kmutnbBlue},
  commentstyle    = \itshape\color{gray},
  stringstyle     = \color{tipteal},
  numberstyle     = \tiny\color{midGray},
  numbers         = left,
  numbersep       = 6pt,
  stepnumber      = 1,
  backgroundcolor = \color{codeGray},
  frame           = single,
  framerule       = 0.4pt,
  rulecolor       = \color{codeFrame},
  breaklines      = true,
  showstringspaces= false,
  tabsize         = 4,
  xleftmargin     = 1.5em,
  xrightmargin    = 0.5em,
  aboveskip       = 0.8em,
  belowskip       = 0.8em,
}

\lstdefinestyle{bashstyle}{
  language        = bash,
  basicstyle      = \ttfamily\small,
  keywordstyle    = \color{kmutnbBlue},
  commentstyle    = \itshape\color{gray},
  backgroundcolor = \color{black!5},
  frame           = single,
  framerule       = 0.4pt,
  rulecolor       = \color{codeFrame},
  breaklines      = true,
  showstringspaces= false,
  xleftmargin     = 1.5em,
  xrightmargin    = 0.5em,
  aboveskip       = 0.6em,
  belowskip       = 0.6em,
}

\lstdefinestyle{textstyle}{
  basicstyle      = \ttfamily\footnotesize,
  backgroundcolor = \color{black!4},
  frame           = single,
  framerule       = 0.4pt,
  rulecolor       = \color{codeFrame},
  breaklines      = true,
  xleftmargin     = 1.5em,
  xrightmargin    = 0.5em,
  aboveskip       = 0.6em,
  belowskip       = 0.6em,
}

% ── Callout boxes ──────────────────────────────────────────────────────────────
\tcbuselibrary{skins, breakable}

\newtcolorbox{tipbox}[1][]{
  enhanced, breakable,
  colback=tipteal!8, colframe=tipteal, coltitle=white,
  fonttitle=\bfseries\small, title={\faLightbulb\quad Tip#1},
  left=6pt, right=6pt, top=4pt, bottom=4pt,
  boxrule=0.8pt, arc=3pt,
}

\newtcolorbox{warnbox}[1][]{
  enhanced, breakable,
  colback=warnAmber!8, colframe=kmutnbGold, coltitle=white,
  fonttitle=\bfseries\small, title={\faExclamationTriangle\quad Note#1},
  left=6pt, right=6pt, top=4pt, bottom=4pt,
  boxrule=0.8pt, arc=3pt,
}

\newtcolorbox{checkpointbox}[1][]{
  enhanced,
  colback=kmutnbBlue!6, colframe=kmutnbBlue, coltitle=white,
  fonttitle=\bfseries\small, title={\faCheckCircle\quad Checkpoint#1},
  left=6pt, right=6pt, top=4pt, bottom=4pt,
  boxrule=0.8pt, arc=3pt,
}

\newtcolorbox{questionbox}[1][]{
  enhanced, breakable,
  colback=kmutnbGold!10, colframe=kmutnbGold, coltitle=white,
  fonttitle=\bfseries\small, title={\faQuestion\quad Question#1},
  left=6pt, right=6pt, top=4pt, bottom=4pt,
  boxrule=0.8pt, arc=3pt,
}

% ── Checkbox list ──────────────────────────────────────────────────────────────
\newlist{checklist}{itemize}{1}
\setlist[checklist]{label=\faSquare[regular], leftmargin=1.5em, noitemsep}

% ── Misc helpers ───────────────────────────────────────────────────────────────
\newcommand{\file}[1]{\texttt{\small #1}}
\newcommand{\api}[1]{\texttt{\small #1}}

% ── Title block macro ─────────────────────────────────────────────────────────
% Usage: \labtitle{03}{Aperiodic Servers \& Producer--Consumer}{2}{CLO 1 \& CLO 2}
\newcommand{\labtitle}[4]{%
  \begin{center}
    {\color{kmutnbBlue}\rule{\linewidth}{2pt}}\\[6pt]
    {\LARGE\bfseries\color{kmutnbBlue} Laboratory #1}\\[4pt]
    {\large\bfseries #2}\\[2pt]
    {\normalsize Real-Time Operating Systems \enspace\textbullet\enspace
     M.Eng.\ ECE \enspace\textbullet\enspace KMUTNB}\\[8pt]
    {\color{kmutnbGold}\rule{\linewidth}{1pt}}
  \end{center}
  \vspace{0.3cm}
  \begin{tabular}{@{}p{3.8cm} p{10cm}@{}}
    \textbf{Platform}       & NXP FRDM-MCXN236 (ARM Cortex-M33 @ 150\,MHz) \\[2pt]
    \textbf{RTOS}           & FreeRTOS v10.6 \\[2pt]
    \textbf{Toolchain}      & ARM GNU Toolchain (\texttt{arm-none-eabi-gcc}) + Make \\[2pt]
    \textbf{Estimated time} & #3 hours \\[2pt]
    \textbf{CLO mapping}    & #4 \\
  \end{tabular}
  \vspace{0.4cm}
}
  (from each labNN/ directory)
% ──────────────────────────────────────────────────────────────────────────────

% ── Packages ──────────────────────────────────────────────────────────────────
\usepackage[a4paper, top=2.5cm, bottom=2.5cm, left=2.5cm, right=2.5cm, headheight=28pt]{geometry}
\usepackage[utf8]{inputenc}
\usepackage[T1]{fontenc}
\usepackage{lmodern}
\usepackage{booktabs}
\usepackage{array}
\usepackage{tabularx}
\usepackage{longtable}
\usepackage{multirow}
\usepackage{colortbl}
\usepackage{xcolor}
\usepackage{titlesec}
\usepackage{enumitem}
\usepackage{hyperref}
\usepackage{fancyhdr}
\usepackage{parskip}
\usepackage{listings}
\usepackage[most]{tcolorbox}
\usepackage{fontawesome5}
\usepackage{microtype}
\usepackage{graphicx}
\usepackage{amsmath}

% ── Colors (KMUTNB brand) ──────────────────────────────────────────────────────
\definecolor{kmutnbBlue}{RGB}{0,56,116}
\definecolor{kmutnbGold}{RGB}{186,146,36}
\definecolor{lightGray}{RGB}{245,245,245}
\definecolor{midGray}{RGB}{200,200,200}
\definecolor{codeGray}{RGB}{248,248,248}
\definecolor{codeFrame}{RGB}{220,220,220}
\definecolor{tipteal}{RGB}{0,105,92}
\definecolor{warnAmber}{RGB}{121,85,0}
\definecolor{checkGreen}{RGB}{27,94,32}

% ── Section formatting ─────────────────────────────────────────────────────────
\titleformat{\section}
  {\large\bfseries\color{kmutnbBlue}}
  {\thesection.}{0.5em}{}
  [\color{kmutnbGold}\titlerule]

\titleformat{\subsection}
  {\normalsize\bfseries\color{kmutnbBlue}}
  {\thesubsection.}{0.5em}{}

\titleformat{\subsubsection}
  {\small\bfseries\color{kmutnbBlue}}
  {\thesubsubsection.}{0.5em}{}

\titlespacing*{\section}{0pt}{1.4ex plus .2ex}{0.6ex plus .1ex}
\titlespacing*{\subsection}{0pt}{1.0ex plus .2ex}{0.4ex plus .1ex}

% ── Listings (C and bash) ──────────────────────────────────────────────────────
\lstdefinestyle{cstyle}{
  language        = C,
  basicstyle      = \ttfamily\small,
  keywordstyle    = \bfseries\color{kmutnbBlue},
  commentstyle    = \itshape\color{gray},
  stringstyle     = \color{tipteal},
  numberstyle     = \tiny\color{midGray},
  numbers         = left,
  numbersep       = 6pt,
  stepnumber      = 1,
  backgroundcolor = \color{codeGray},
  frame           = single,
  framerule       = 0.4pt,
  rulecolor       = \color{codeFrame},
  breaklines      = true,
  showstringspaces= false,
  tabsize         = 4,
  xleftmargin     = 1.5em,
  xrightmargin    = 0.5em,
  aboveskip       = 0.8em,
  belowskip       = 0.8em,
}

\lstdefinestyle{bashstyle}{
  language        = bash,
  basicstyle      = \ttfamily\small,
  keywordstyle    = \color{kmutnbBlue},
  commentstyle    = \itshape\color{gray},
  backgroundcolor = \color{black!5},
  frame           = single,
  framerule       = 0.4pt,
  rulecolor       = \color{codeFrame},
  breaklines      = true,
  showstringspaces= false,
  xleftmargin     = 1.5em,
  xrightmargin    = 0.5em,
  aboveskip       = 0.6em,
  belowskip       = 0.6em,
}

\lstdefinestyle{textstyle}{
  basicstyle      = \ttfamily\footnotesize,
  backgroundcolor = \color{black!4},
  frame           = single,
  framerule       = 0.4pt,
  rulecolor       = \color{codeFrame},
  breaklines      = true,
  xleftmargin     = 1.5em,
  xrightmargin    = 0.5em,
  aboveskip       = 0.6em,
  belowskip       = 0.6em,
}

% ── Callout boxes ──────────────────────────────────────────────────────────────
\tcbuselibrary{skins, breakable}

\newtcolorbox{tipbox}[1][]{
  enhanced, breakable,
  colback=tipteal!8, colframe=tipteal, coltitle=white,
  fonttitle=\bfseries\small, title={\faLightbulb\quad Tip#1},
  left=6pt, right=6pt, top=4pt, bottom=4pt,
  boxrule=0.8pt, arc=3pt,
}

\newtcolorbox{warnbox}[1][]{
  enhanced, breakable,
  colback=warnAmber!8, colframe=kmutnbGold, coltitle=white,
  fonttitle=\bfseries\small, title={\faExclamationTriangle\quad Note#1},
  left=6pt, right=6pt, top=4pt, bottom=4pt,
  boxrule=0.8pt, arc=3pt,
}

\newtcolorbox{checkpointbox}[1][]{
  enhanced,
  colback=kmutnbBlue!6, colframe=kmutnbBlue, coltitle=white,
  fonttitle=\bfseries\small, title={\faCheckCircle\quad Checkpoint#1},
  left=6pt, right=6pt, top=4pt, bottom=4pt,
  boxrule=0.8pt, arc=3pt,
}

\newtcolorbox{questionbox}[1][]{
  enhanced, breakable,
  colback=kmutnbGold!10, colframe=kmutnbGold, coltitle=white,
  fonttitle=\bfseries\small, title={\faQuestion\quad Question#1},
  left=6pt, right=6pt, top=4pt, bottom=4pt,
  boxrule=0.8pt, arc=3pt,
}

% ── Checkbox list ──────────────────────────────────────────────────────────────
\newlist{checklist}{itemize}{1}
\setlist[checklist]{label=\faSquare[regular], leftmargin=1.5em, noitemsep}

% ── Misc helpers ───────────────────────────────────────────────────────────────
\newcommand{\file}[1]{\texttt{\small #1}}
\newcommand{\api}[1]{\texttt{\small #1}}

% ── Title block macro ─────────────────────────────────────────────────────────
% Usage: \labtitle{03}{Aperiodic Servers \& Producer--Consumer}{2}{CLO 1 \& CLO 2}
\newcommand{\labtitle}[4]{%
  \begin{center}
    {\color{kmutnbBlue}\rule{\linewidth}{2pt}}\\[6pt]
    {\LARGE\bfseries\color{kmutnbBlue} Laboratory #1}\\[4pt]
    {\large\bfseries #2}\\[2pt]
    {\normalsize Real-Time Operating Systems \enspace\textbullet\enspace
     M.Eng.\ ECE \enspace\textbullet\enspace KMUTNB}\\[8pt]
    {\color{kmutnbGold}\rule{\linewidth}{1pt}}
  \end{center}
  \vspace{0.3cm}
  \begin{tabular}{@{}p{3.8cm} p{10cm}@{}}
    \textbf{Platform}       & NXP FRDM-MCXN236 (ARM Cortex-M33 @ 150\,MHz) \\[2pt]
    \textbf{RTOS}           & FreeRTOS v10.6 \\[2pt]
    \textbf{Toolchain}      & ARM GNU Toolchain (\texttt{arm-none-eabi-gcc}) + Make \\[2pt]
    \textbf{Estimated time} & #3 hours \\[2pt]
    \textbf{CLO mapping}    & #4 \\
  \end{tabular}
  \vspace{0.4cm}
}
  (from each labNN/ directory)
% ──────────────────────────────────────────────────────────────────────────────

% ── Packages ──────────────────────────────────────────────────────────────────
\usepackage[a4paper, top=2.5cm, bottom=2.5cm, left=2.5cm, right=2.5cm, headheight=28pt]{geometry}
\usepackage[utf8]{inputenc}
\usepackage[T1]{fontenc}
\usepackage{lmodern}
\usepackage{booktabs}
\usepackage{array}
\usepackage{tabularx}
\usepackage{longtable}
\usepackage{multirow}
\usepackage{colortbl}
\usepackage{xcolor}
\usepackage{titlesec}
\usepackage{enumitem}
\usepackage{hyperref}
\usepackage{fancyhdr}
\usepackage{parskip}
\usepackage{listings}
\usepackage[most]{tcolorbox}
\usepackage{fontawesome5}
\usepackage{microtype}
\usepackage{graphicx}
\usepackage{amsmath}

% ── Colors (KMUTNB brand) ──────────────────────────────────────────────────────
\definecolor{kmutnbBlue}{RGB}{0,56,116}
\definecolor{kmutnbGold}{RGB}{186,146,36}
\definecolor{lightGray}{RGB}{245,245,245}
\definecolor{midGray}{RGB}{200,200,200}
\definecolor{codeGray}{RGB}{248,248,248}
\definecolor{codeFrame}{RGB}{220,220,220}
\definecolor{tipteal}{RGB}{0,105,92}
\definecolor{warnAmber}{RGB}{121,85,0}
\definecolor{checkGreen}{RGB}{27,94,32}

% ── Section formatting ─────────────────────────────────────────────────────────
\titleformat{\section}
  {\large\bfseries\color{kmutnbBlue}}
  {\thesection.}{0.5em}{}
  [\color{kmutnbGold}\titlerule]

\titleformat{\subsection}
  {\normalsize\bfseries\color{kmutnbBlue}}
  {\thesubsection.}{0.5em}{}

\titleformat{\subsubsection}
  {\small\bfseries\color{kmutnbBlue}}
  {\thesubsubsection.}{0.5em}{}

\titlespacing*{\section}{0pt}{1.4ex plus .2ex}{0.6ex plus .1ex}
\titlespacing*{\subsection}{0pt}{1.0ex plus .2ex}{0.4ex plus .1ex}

% ── Listings (C and bash) ──────────────────────────────────────────────────────
\lstdefinestyle{cstyle}{
  language        = C,
  basicstyle      = \ttfamily\small,
  keywordstyle    = \bfseries\color{kmutnbBlue},
  commentstyle    = \itshape\color{gray},
  stringstyle     = \color{tipteal},
  numberstyle     = \tiny\color{midGray},
  numbers         = left,
  numbersep       = 6pt,
  stepnumber      = 1,
  backgroundcolor = \color{codeGray},
  frame           = single,
  framerule       = 0.4pt,
  rulecolor       = \color{codeFrame},
  breaklines      = true,
  showstringspaces= false,
  tabsize         = 4,
  xleftmargin     = 1.5em,
  xrightmargin    = 0.5em,
  aboveskip       = 0.8em,
  belowskip       = 0.8em,
}

\lstdefinestyle{bashstyle}{
  language        = bash,
  basicstyle      = \ttfamily\small,
  keywordstyle    = \color{kmutnbBlue},
  commentstyle    = \itshape\color{gray},
  backgroundcolor = \color{black!5},
  frame           = single,
  framerule       = 0.4pt,
  rulecolor       = \color{codeFrame},
  breaklines      = true,
  showstringspaces= false,
  xleftmargin     = 1.5em,
  xrightmargin    = 0.5em,
  aboveskip       = 0.6em,
  belowskip       = 0.6em,
}

\lstdefinestyle{textstyle}{
  basicstyle      = \ttfamily\footnotesize,
  backgroundcolor = \color{black!4},
  frame           = single,
  framerule       = 0.4pt,
  rulecolor       = \color{codeFrame},
  breaklines      = true,
  xleftmargin     = 1.5em,
  xrightmargin    = 0.5em,
  aboveskip       = 0.6em,
  belowskip       = 0.6em,
}

% ── Callout boxes ──────────────────────────────────────────────────────────────
\tcbuselibrary{skins, breakable}

\newtcolorbox{tipbox}[1][]{
  enhanced, breakable,
  colback=tipteal!8, colframe=tipteal, coltitle=white,
  fonttitle=\bfseries\small, title={\faLightbulb\quad Tip#1},
  left=6pt, right=6pt, top=4pt, bottom=4pt,
  boxrule=0.8pt, arc=3pt,
}

\newtcolorbox{warnbox}[1][]{
  enhanced, breakable,
  colback=warnAmber!8, colframe=kmutnbGold, coltitle=white,
  fonttitle=\bfseries\small, title={\faExclamationTriangle\quad Note#1},
  left=6pt, right=6pt, top=4pt, bottom=4pt,
  boxrule=0.8pt, arc=3pt,
}

\newtcolorbox{checkpointbox}[1][]{
  enhanced,
  colback=kmutnbBlue!6, colframe=kmutnbBlue, coltitle=white,
  fonttitle=\bfseries\small, title={\faCheckCircle\quad Checkpoint#1},
  left=6pt, right=6pt, top=4pt, bottom=4pt,
  boxrule=0.8pt, arc=3pt,
}

\newtcolorbox{questionbox}[1][]{
  enhanced, breakable,
  colback=kmutnbGold!10, colframe=kmutnbGold, coltitle=white,
  fonttitle=\bfseries\small, title={\faQuestion\quad Question#1},
  left=6pt, right=6pt, top=4pt, bottom=4pt,
  boxrule=0.8pt, arc=3pt,
}

% ── Checkbox list ──────────────────────────────────────────────────────────────
\newlist{checklist}{itemize}{1}
\setlist[checklist]{label=\faSquare[regular], leftmargin=1.5em, noitemsep}

% ── Misc helpers ───────────────────────────────────────────────────────────────
\newcommand{\file}[1]{\texttt{\small #1}}
\newcommand{\api}[1]{\texttt{\small #1}}

% ── Title block macro ─────────────────────────────────────────────────────────
% Usage: \labtitle{03}{Aperiodic Servers \& Producer--Consumer}{2}{CLO 1 \& CLO 2}
\newcommand{\labtitle}[4]{%
  \begin{center}
    {\color{kmutnbBlue}\rule{\linewidth}{2pt}}\\[6pt]
    {\LARGE\bfseries\color{kmutnbBlue} Laboratory #1}\\[4pt]
    {\large\bfseries #2}\\[2pt]
    {\normalsize Real-Time Operating Systems \enspace\textbullet\enspace
     M.Eng.\ ECE \enspace\textbullet\enspace KMUTNB}\\[8pt]
    {\color{kmutnbGold}\rule{\linewidth}{1pt}}
  \end{center}
  \vspace{0.3cm}
  \begin{tabular}{@{}p{3.8cm} p{10cm}@{}}
    \textbf{Platform}       & NXP FRDM-MCXN236 (ARM Cortex-M33 @ 150\,MHz) \\[2pt]
    \textbf{RTOS}           & FreeRTOS v10.6 \\[2pt]
    \textbf{Toolchain}      & ARM GNU Toolchain (\texttt{arm-none-eabi-gcc}) + Make \\[2pt]
    \textbf{Estimated time} & #3 hours \\[2pt]
    \textbf{CLO mapping}    & #4 \\
  \end{tabular}
  \vspace{0.4cm}
}


% ── Header / Footer ────────────────────────────────────────────────────────────
\pagestyle{fancy}
\fancyhf{}
\fancyhead[L]{\small\color{kmutnbBlue}\textbf{KMUTNB -- Faculty of Engineering}}
\fancyhead[R]{\small\color{kmutnbBlue}Dept.\ of Electrical and Computer Engineering}
\fancyfoot[L]{\small\color{kmutnbBlue}Real-Time Operating Systems -- M.Eng.\ ECE}
\fancyfoot[C]{\small\thepage}
\fancyfoot[R]{\small\color{kmutnbBlue}Lab 10}
\renewcommand{\headrulewidth}{0.4pt}
\renewcommand{\headrule}{\hbox to\headwidth{%
  \color{kmutnbGold}\leaders\hrule height\headrulewidth\hfill}}
\renewcommand{\footrulewidth}{0.4pt}
\renewcommand{\footrule}{\hbox to\headwidth{%
  \color{midGray}\leaders\hrule height\footrulewidth\hfill}}

\hypersetup{
  colorlinks = true,
  linkcolor  = kmutnbBlue,
  urlcolor   = kmutnbBlue,
  citecolor  = kmutnbBlue,
  pdftitle   = {Lab 10 -- Tickless Idle and Power Measurement},
  pdfauthor  = {KMUTNB RTOS Course},
}

% ══════════════════════════════════════════════════════════════════════════════
\begin{document}

\labtitle{10}{Tickless Idle \& Power Measurement}{3--4}{CLO 4 \& CLO 5}

% ── Objectives ────────────────────────────────────────────────────────────────
\section{Objectives}

By completing this lab you will be able to:

\begin{enumerate}[noitemsep, leftmargin=*]
  \item \textbf{Measure} average current consumption with and without FreeRTOS tickless
        idle.
  \item \textbf{Configure} the built-in tickless idle
        (\texttt{configUSE\_TICKLESS\_IDLE\,=\,1}).
  \item \textbf{Implement} a custom \api{vPortSuppressTicksAndSleep} using the MCXN236
        LPTMR.
  \item \textbf{Verify} tick correction accuracy using a 1-second LED toggle and DWT
        timing.
  \item \textbf{Estimate} battery life improvement from tickless idle over a 24-hour
        scenario.
\end{enumerate}

% ── Prerequisites ─────────────────────────────────────────────────────────────
\section{Prerequisites}

\begin{checklist}
  \item Lab 06 complete (DWT cycle counter working).
  \item A current probe or NXP MCU-Link power measurement capability \textbf{OR} an
        ammeter in series with the 3.3\,V supply rail.
  \item ARM GNU Toolchain and \texttt{pyocd} working.
\end{checklist}

\begin{warnbox}
  \textbf{Power measurement alternatives:} If no current probe is available, you can
  use an MCU-Link Pro board's power measurement feature via MCUXpresso IDE, or a USB
  power meter on the host port. The relative comparison (with vs.\ without tickless)
  is more important than the absolute value.
\end{warnbox}

% ── Background ────────────────────────────────────────────────────────────────
\section{Background}

\subsection{The Tick Interrupt Problem}

FreeRTOS generates a SysTick interrupt every 1\,ms (at 1\,kHz tick rate). Even when all
tasks are blocked, the CPU wakes every 1\,ms to run the tick handler --- about 100--200
cycles of overhead, repeated 1\,000 times per second.

On a system where the active task runs only 2\,ms every 500\,ms (0.4\% duty cycle), the
tick alone dominates power consumption in the idle period.

\subsection{Tickless Idle Flow}

\begin{lstlisting}[style=textstyle]
Scheduler finds all tasks blocked:
  1. Compute xExpectedIdleTime = ticks until next task unblocks
  2. Call vPortSuppressTicksAndSleep(xExpectedIdleTime)
     a. Disable SysTick
     b. Program LPTMR for xExpectedIdleTime (in 32 kHz ticks)
     c. __WFI with SLEEPDEEP
     d. Wake (LPTMR or other interrupt)
     e. Read elapsed LPTMR ticks
     f. vTaskStepTick(elapsed) -- correct the tick count
     g. Re-enable SysTick
  3. Resume normal scheduling
\end{lstlisting}

% ── Project Structure ─────────────────────────────────────────────────────────
\section{Project Structure}

\begin{lstlisting}[style=textstyle]
lab10_tickless/
  Makefile
  linker_MCXN236.ld
  src/
    main.c                    <- task creation, LED init, DWT init
    adc_task.c/.h             <- ADC sample task (500 ms period)
    led_toggle_task.c/.h      <- 1-second LED (timing accuracy test)
    tickless_port.c/.h        <- Part C: custom vPortSuppressTicksAndSleep
    FreeRTOSConfig.h
    uart.h / uart.c
  FreeRTOS-Kernel/
\end{lstlisting}

% ── Part A ─────────────────────────────────────────────────────────────────────
\section{Part A --- Baseline Current Measurement (No Tickless)}

\subsection{Step 1 --- Configure with Tickless Disabled}

In \file{FreeRTOSConfig.h}:

\begin{lstlisting}[style=cstyle]
#define configUSE_TICKLESS_IDLE         0
#define configTICK_RATE_HZ              1000
\end{lstlisting}

\subsection{Step 2 --- Implement the ADC Task}

\begin{lstlisting}[style=cstyle]
void vADCTask(void *pv)
{
    TickType_t xLastWake = xTaskGetTickCount();
    uint32_t sample = 0;
    for (;;) {
        vTaskDelayUntil(&xLastWake, pdMS_TO_TICKS(500));
        /* Simulate ADC read: ~50 cycles */
        sample = (sample + 37) % 4096;
        uart_printf("[ADC] sample=%lu\r\n", sample);
    }
}
\end{lstlisting}

This task is active for $\sim$50 cycles (0.33\,\textmu s) every 500\,ms --- duty cycle
$\approx$ 0.000067\%.

\subsection{Step 3 --- Measure Current}

Connect your ammeter or current probe in series with the board's 3.3\,V supply. Run
for 30 seconds and record minimum and average current.

\begin{checkpointbox}[~A]
  Record average mA with \texttt{configUSE\_TICKLESS\_IDLE\,=\,0}.
\end{checkpointbox}

\begin{questionbox}[~A1]
  The tick fires 1\,000 times per second. Each tick handler takes roughly 150 cycles.
  What fraction of total CPU time is consumed by tick handling alone? Convert to a
  percentage.
\end{questionbox}

% ── Part B ─────────────────────────────────────────────────────────────────────
\section{Part B --- Built-In Tickless Idle}

\subsection{Step 1 --- Enable Built-In Tickless}

In \file{FreeRTOSConfig.h}:

\begin{lstlisting}[style=cstyle]
#define configUSE_TICKLESS_IDLE                  1
#define configEXPECTED_IDLE_TIME_BEFORE_SLEEP    2
\end{lstlisting}

No other code changes are needed --- the Cortex-M33 FreeRTOS port implements
\api{vPortSuppressTicksAndSleep} using SysTick reload.

\subsection{Step 2 --- Measure Current Again}

Repeat the same 30-second measurement.

\begin{checkpointbox}[~B]
  Current measurement with \texttt{configUSE\_TICKLESS\_IDLE\,=\,1}.
\end{checkpointbox}

\begin{questionbox}[~B1]
  How much did the average current drop? Express as a percentage reduction.
\end{questionbox}

\begin{questionbox}[~B2]
  With tickless enabled, the CPU still wakes every 500\,ms for the ADC task.
  Calculate the total number of CPU wakeups per hour:
  \begin{itemize}[noitemsep]
    \item Without tickless: \underline{\hspace{2cm}} (tick + task)
    \item With tickless:    \underline{\hspace{2cm}} (task only)
  \end{itemize}
\end{questionbox}

% ── Part C ─────────────────────────────────────────────────────────────────────
\section{Part C --- Custom LPTMR Tickless Implementation}

The built-in port uses SysTick (which stops in Deep Sleep on some MCU configurations).
A better approach is to use the LPTMR, which runs from the 32.768\,kHz VBAT clock and
survives Deep Sleep.

\subsection{Step 1 --- Switch to Custom Tickless}

In \file{FreeRTOSConfig.h}:

\begin{lstlisting}[style=cstyle]
#define configUSE_TICKLESS_IDLE         2   /* custom implementation */
\end{lstlisting}

\subsection{Step 2 --- Implement \texttt{vPortSuppressTicksAndSleep}}

In \file{tickless\_port.c}, implement all six steps:

\begin{lstlisting}[style=cstyle]
#include "FreeRTOS.h"
#include "task.h"
#include "fsl_lptmr.h"

#define LPTMR_CLOCK_HZ   32768UL   /* VBAT 32 kHz oscillator */

void vPortSuppressTicksAndSleep(TickType_t xExpectedIdleTime)
{
    /* --- 1. Disable SysTick --- */
    SysTick->CTRL &= ~SysTick_CTRL_ENABLE_Msk;

    /* --- 2. Program LPTMR --- */
    uint32_t lptmr_ticks = (uint32_t)(
        ((uint64_t)xExpectedIdleTime * LPTMR_CLOCK_HZ)
        / (uint64_t)configTICK_RATE_HZ);

    if (lptmr_ticks == 0) {
        SysTick->CTRL |= SysTick_CTRL_ENABLE_Msk;
        return;
    }

    LPTMR0->CSR = 0;
    LPTMR0->CMR = lptmr_ticks - 1;
    LPTMR0->CSR = LPTMR_CSR_TEN_MASK | LPTMR_CSR_TIE_MASK;

    /* --- 3. Enter Deep Sleep --- */
    SCB->SCR |= SCB_SCR_SLEEPDEEP_Msk;
    __DSB();
    __WFI();
    SCB->SCR &= ~SCB_SCR_SLEEPDEEP_Msk;

    /* --- 4. Stop LPTMR and read elapsed time --- */
    LPTMR0->CSR &= ~LPTMR_CSR_TEN_MASK;
    uint32_t elapsed_lptmr = LPTMR0->CNR;
    LPTMR0->CSR |= LPTMR_CSR_TCF_MASK;

    /* --- 5. Correct FreeRTOS tick count --- */
    TickType_t elapsed_ticks = (TickType_t)(
        ((uint64_t)elapsed_lptmr * configTICK_RATE_HZ)
        / (uint64_t)LPTMR_CLOCK_HZ);

    if (elapsed_ticks > 0)
        vTaskStepTick(elapsed_ticks);

    /* --- 6. Re-enable SysTick --- */
    SysTick->CTRL |= SysTick_CTRL_ENABLE_Msk;
}

void LPTMR0_IRQHandler(void)
{
    LPTMR0->CSR |= LPTMR_CSR_TCF_MASK;   /* clear flag -- wake only */
}
\end{lstlisting}

\subsection{Step 3 --- Verify Timing Accuracy with LED Toggle}

\begin{lstlisting}[style=cstyle]
void vLEDTask(void *pv)
{
    TickType_t xLastWake = xTaskGetTickCount();
    uint32_t toggle_count = 0;
    DWT_Enable();
    uint32_t prev_cycles = DWT->CYCCNT;

    for (;;) {
        vTaskDelayUntil(&xLastWake, pdMS_TO_TICKS(1000));
        uint32_t now_cycles = DWT->CYCCNT;
        uint32_t elapsed_ms = (now_cycles - prev_cycles) / 150000;
        prev_cycles = now_cycles;

        GPIO_TogglePinsOutput(BOARD_LED_GPIO, 1u << BOARD_LED_PIN);
        toggle_count++;
        uart_printf("[LED] toggle #%lu  actual period = %lu ms\r\n",
                    toggle_count, elapsed_ms);
    }
}
\end{lstlisting}

Run for 60 seconds and observe whether the LED period drifts from 1\,000\,ms.

\begin{checkpointbox}[~C]
  UART from the LED task showing 60 toggle events with period values (must stay
  within 1\,000 $\pm$ 2\,ms).
\end{checkpointbox}

\begin{questionbox}[~C1]
  The LPTMR runs at 32.768\,kHz. The tick is 1\,kHz. What is the maximum rounding
  error when converting between LPTMR ticks and FreeRTOS ticks? Express as $\pm$\,ms.
\end{questionbox}

\begin{questionbox}[~C2]
  \api{vTaskStepTick} must be called with a value \textbf{less than}
  \texttt{xExpectedIdleTime}. Why? What happens if you pass a value that is too large?
  (Hint: look at \file{task.c} --- \api{vTaskStepTick} asserts.)
\end{questionbox}

% ── Part D ─────────────────────────────────────────────────────────────────────
\section{Part D --- Energy Analysis}

\subsection{Step 1 --- Measure Current with Custom LPTMR Implementation}

Repeat the 30-second current measurement from Parts A and B.

\subsection{Step 2 --- Complete the Comparison Table}

\renewcommand{\arraystretch}{1.4}
\begin{tabular}{@{} p{4.5cm} p{2.8cm} p{2.4cm} p{2.8cm} @{}}
  \toprule
  \textbf{Configuration} & \textbf{Avg.\ current (mA)} & \textbf{Wakeups/hour} &
  \textbf{Energy (mWh/24h)} \\
  \midrule
  No tickless (Part A) & & 3\,600\,000 & \\
  Built-in tickless (Part B) & & 7\,200 & \\
  Custom LPTMR (Part D) & & 7\,200 & \\
  \bottomrule
\end{tabular}
\renewcommand{\arraystretch}{1}

\medskip
Energy (mWh) = average\_current (mA) $\times$ 3.3\,V $\times$ 24\,h $\div$ 1000.

\begin{questionbox}[~D1]
  Using your measured currents, how many days would a CR2032 coin cell (225\,mAh
  capacity) power the board under each configuration?
\end{questionbox}

\begin{questionbox}[~D2]
  \api{vTaskDelayUntil} in \texttt{vADCTask} wakes the task at exactly 500\,ms
  boundaries. If the LPTMR rounding error from Question C1 is $\pm$0.5\,ms per sleep
  cycle, what is the maximum accumulated timing error after 1 hour? Is this acceptable
  for a real-time sensor system?
\end{questionbox}

% ── Deliverables ───────────────────────────────────────────────────────────────
\section{Deliverables}

\renewcommand{\arraystretch}{1.3}
\begin{tabular}{@{} c p{5cm} p{6.5cm} @{}}
  \toprule
  \textbf{\#} & \textbf{Item} & \textbf{Details} \\
  \midrule
  1 & Part A current & Value and methodology description \\
  2 & Part B current & Value (tickless = 1) \\
  3 & Part C UART & 60 LED toggle events, period within 1\,000 $\pm$ 2\,ms \\
  4 & Part D table & Completed comparison with measured values \\
  5 & \file{tickless\_port.c} & Source listing \\
  6 & Written answers & Questions A1, B1--B2, C1--C2, D1--D2 \\
  \bottomrule
\end{tabular}
\renewcommand{\arraystretch}{1}

% ── Key Concepts ───────────────────────────────────────────────────────────────
\section{Key Concepts Demonstrated}

\renewcommand{\arraystretch}{1.3}
\begin{tabular}{@{} p{6.5cm} p{6cm} @{}}
  \toprule
  \textbf{Concept} & \textbf{Where you saw it} \\
  \midrule
  Tick interrupt power overhead & Part A current measurement \\
  Built-in SysTick tickless idle & Part B \\
  LPTMR-based custom tickless port & Part C \\
  \api{vTaskStepTick} tick correction & Part C \\
  Timing accuracy with LPTMR rounding & Part C LED toggle \\
  Energy analysis and battery life estimation & Part D \\
  \bottomrule
\end{tabular}
\renewcommand{\arraystretch}{1}

% ── Reference ─────────────────────────────────────────────────────────────────
\section{Reference}

\renewcommand{\arraystretch}{1.3}
\begin{tabular}{@{} p{5.5cm} p{7cm} @{}}
  \toprule
  \textbf{Resource} & \textbf{Relevant sections} \\
  \midrule
  FreeRTOS docs & \api{configUSE\_TICKLESS\_IDLE}, \api{vTaskStepTick},
                  \api{vPortSuppressTicksAndSleep} \\
  NXP MCXN236 Reference Manual & Chapter PMC, Chapter LPTMR \\
  ARM Cortex-M33 TRM & \S B2.5 Power management, WFI instruction \\
  NXP MCUXpresso SDK & \file{fsl\_lptmr.h} \\
  \bottomrule
\end{tabular}
\renewcommand{\arraystretch}{1}

\end{document}
